\section{Conclusion}
\label{sec:conclusion}

In this paper, we introduced \textbf{Cross-Attention Multi-Expert Routing (CAM-Router)}, a lightweight and highly robust dynamic model merging framework designed to resolve parameter conflicts under orthogonal task configurations. By rejecting the standard practice of immediate global average pooling, CAM-Router preserves the full spatial token sequence from the early layers of the transformer backbone. Trainable task-expert queries learn to dynamically attend to localized spatial regions via Multi-Head Cross-Attention (MHCA), capturing specialized domain features that are lost in global pooling. Furthermore, we implemented independent bounded sigmoidal gating, eliminating the zero-sum competitive pressure imposed by traditional Softmax constraints.

Guided by our \textbf{Empiricist} research methodology, we subjected CAM-Router to rigorous, large-scale empirical evaluation. Our experimental sandbox—simulated on a compact 14-layer Vision Transformer backbone across MNIST, FashionMNIST, CIFAR-10, and SVHN—verified that CAM-Router achieves a breakthrough Joint Mean Accuracy of \textbf{53.07\%}. This represents a massive \textbf{+11.10\%} absolute accuracy improvement over the Static Uniform baseline (41.97\%) and significantly outperforms both classical and quantum average-pooling routers. In addition, systematic stress testing demonstrated that CAM-Router exhibits exceptional resilience to spatial occlusions (maintaining a stable accuracy of \textbf{50.57\%} under 80\% patch masking) and is immune to the "heterogeneity collapse" that plagues other average-pooling methods under large mixed-task batches ($B=256$, achieving \textbf{54.30\%}).

We hope our findings encourage the model-merging community to move away from rigid global pooling abstractions and embrace spatially aware, empirically validated weight routing structures. Interesting future directions include scaling CAM-Router to larger foundation architectures (e.g., LLaMA and multi-modal CLIP backbones) and investigating the potential of token-level dynamic parameter fusions.
